\documentclass[slidestop,serif,compress]{beamer}    %slidestop=title in up-left corner (with "slidescentered" centered); compress=navigation bars as small as possible
\usepackage[T1]{fontenc}
\usepackage{lmodern}
\usepackage[ugm]{mathdesign} %use Garamond font

%gets rid of bottom navigation bars and puts slide numbers right bottom
\setbeamertemplate{footline}[frame number]{}

%alternative footer: name, institute, page number
%\setbeamertemplate{footline}
%{%
%  \leavevmode%
%  \hbox{\begin{beamercolorbox}[wd=.5\paperwidth,ht=2.5ex,dp=1.125ex,leftskip=.3cm plus1fill,rightskip=.3cm]{author in head/foot}%
%    \usebeamerfont{author in head/foot}\insertshortauthor
%  \end{beamercolorbox}%
%  \begin{beamercolorbox}[wd=.5\paperwidth,ht=2.5ex,dp=1.125ex,leftskip=.3cm,rightskip=.3cm plus1fil]{title in head/foot}%
%    \usebeamerfont{title in head/foot}\insertshortinstitute \hfill \insertframenumber \,/\,\inserttotalframenumber
%  \end{beamercolorbox}}%
%  \vskip0pt%
%}

%gets rid of navigation symbols
\setbeamertemplate{navigation symbols}{}


%%complete themes (not recommended)
%%\usetheme{Madrid}             %"Berkeley":lhs navigation bar; 1st alternative: "Madrid" without a navigation bar but with a footer including name (institution), title, date page numbering (but then \institute{Univ. ...}); 2nd alternative: defaults without all the extras; 3rd alternative "Dresden": two bottom lines)
%%"Dresden is possibly easy to adapt to UvT colors
%instead of using a complete theme inner and outer themes can be specified in isolation
%inner theme:
\useinnertheme{rounded} %actually not necessary but nice for theorem boxes etc. but enumerate listings look a bit stupid! Alternative: default = just comment it out
%outer themes:
%\useoutertheme[subsection=false,footline=authorinstitutetitlepage]{miniframes} %other options for footer see beamer user guide p. 153


%%complete color themes: (ok but inner outer styles below are better)
%%\usecolortheme{albatross}    % background is colored (in blue) in albatross; the option [overlystylish] installs a background canva which is definitely too stylish for an economics department
%%\usecolortheme{seagull}   %for greyscale (black and white)
%%default is not bad!!!

%alternative to complete color schemes one can define an inner and outer color scheme
%inner color themes
%\usecolortheme{orchid} %deep color in box(theorem etc.) background fits with outer color scheme "whale"
\usecolortheme{rose} % less aggressive slight grey shading background (in boxes/theorems etc.) fits nicely with outer color scheme \usecolortheme{seahorse}
%outer color theme
\usecolortheme{dolphin}   %alternative to seahorse: "whale" but only if inner scheme is "orchid"
%
%\usefonttheme{structurebold}       %titles, headlines, foodlines and sidebars are bold

\usepackage{amsmath, eurosym, textcomp, graphicx, amsthm, amssymb, hyperref,tikz,xcolor,sgamevar}

\setlength {\parindent}{0cm}
\newcommand {\bi}{\begin{itemize}}
\newcommand {\ei}{\end{itemize}}
\newcommand{\be}{\begin{enumerate}}
\newcommand{\ee}{\end{enumerate}}
\newcommand{\Ra}{\Rightarrow}
\newcommand{\ra}{\rightarrow}
\newcommand{\Lra}{\Leftrightarrow}
\newcommand{\del}{\partial}
\newcommand{\bea}{\begin{eqnarray*}}
\newcommand{\eea}{\end{eqnarray*}}

\renewcommand{\Re}{\mathbb{R}}

\theoremstyle{plain}
\newtheorem{prop}{Proposition}
%\newtheorem{theorem}{Theorem}
\newtheorem{assumption}{Assumption}
%\newtheorem{lemma}{Lemma}
\newtheorem{proposition}{Proposition}
\newtheorem{observation}{Observation}
%\newtheorem{corollary}{Corollary}
\theoremstyle{definition}
\newtheorem{ex}{Example}
\newtheorem{exercise}{Exercise}
%\newtheorem{definition}{Definition}

%\setlength{\parskip}{1.5ex plus 0.5ex minus 0.5ex}
\author{Christoph Schottm�ller}
\institute{University of Copenhagen}
\title{Game Theory\\ Nash Theorem\footnote{License: CC Attribution Share Alike 4.0}}
\date{}

\newcommand{\fr}{\frame[allowframebreaks]}
\newcommand{\ft}{\frametitle}

\newcommand{\backupbegin}{
   \newcounter{framenumberappendix}
   \setcounter{framenumberappendix}{\value{framenumber}}
}
\newcommand{\backupend}{
   \addtocounter{framenumberappendix}{-\value{framenumber}}
   \addtocounter{framenumber}{\value{framenumberappendix}}
}

\begin{document}

\begin{frame}   % Cover slide
\titlepage
\end{frame}

\section[Outline]{}
\frame{\ft{Outline}\tableofcontents}

\section{Existence of Nash equilibria}
\fr{\ft{Existence of Nash equilibrium}

  \begin{itemize}
  \item not all games have a Nash equilibrium
  \end{itemize}

  \begin{ex}
    \begin{itemize}
    \item 2 players ($N=\{P1,P2\}$)
    \item each player says a number ($A_i=\Re$)
    \item player that says the higher number wins \\(e.g. winner has payoff 1 while loser has payoff 0, if both say the same number each has payoff 1/2)
    \end{itemize}
  \end{ex}



\framebreak

Why are we interested in existence?
\begin{itemize}
\item worthwhile to search for an equilibrium
\item steady state interpretation of equilibrium, existence of NE says that the process might become stable
\item  sometimes possible to do comparative statics without computing the NE (only sensible if NE exists)
\end{itemize}
}

\section{Maths: Fixed point theorems}

\fr{\ft{Fixed point theory}
  \begin{definition}[fixed point]
    Let $f$ be a real function, i.e. $f:\,A\rightarrow\Re$ where $A\subseteq \Re$. A point $a\in A$ is a \emph{fixed point} of $f$ if $f(a)=a$.
  \end{definition}

 \begin{itemize}
  \item Does $f(x)=1$ have a fixed point? What about $f(x)=x+1$? What about $f(x)=x^2$?
  \item fixed point theory gives general conditions under which functions have fixed points
  \end{itemize}

\framebreak

  \begin{theorem}[Brouwer's fixed point theorem, 1-dimensional]
    Let $f:[0,1]\rightarrow[0,1]$ be a continuous function. Then $f$ has a fixed point.
  \end{theorem}

\framebreak

%Why is Brouwer's theorem interesting for game theorists?
\begin{ex}[2 player Cournot]
  \begin{itemize}
  \item two firms each choose a quantity $q_i$ 
  \item both firms have costs $c(q_i)=c q_i$ for some $c>0$
  \item inverse demand is $P(q_1+q_2)$ where we assume that $P$ is two times continuously differentiable with $P'<0$ and $P''\leq 0$
  \item assume that $P(1)<c$\\ $\Rightarrow$ a firm will never offer a quantity greater than 1
  \item  firm 1 chooses a quantity from $[0,1]$ to maximize profits
   $$\max_{q_1} (P(q_1+q_2)-c)q_1$$
 we get the first order condition
  $$P'(q_1+q_2)q_1+P(q_1+q_2)-c=0$$
  \end{itemize}
\end{ex}

\begin{ex}[2 player Cournot continued]
  \begin{itemize}
  \item the second order condition holds by assumption
$$P''(q_1+q_2)q_1+2P'(q_1+q_2)<0$$
\item the first order condition defines a best response function
  $q_1(q_2)$
\item the best response function is continuous because $P$ and $P'$
  are continuous by assumption
\item Brouwer: best response function has a fixed point!
\item game is symmetric $\Rightarrow$ fixed point is an
  equilibrium $\Rightarrow$ equilibrium exists
\item We showed this without being able to actually calculate the equilibrium!!!
\end{itemize}
\end{ex}
\framebreak

\begin{theorem}[Brouwer fixed point theorem]
  Let $S$ be a convex and compact set in $\Re^n$ and let $f:S\rightarrow S$ be a continuous function. Then there exists an $s^*\in S$ such that $f(s^*)=s^*$.
\end{theorem}
}





\section{Nash theorem}

\frame{\ft{Nash theorem}
\begin{theorem}[Nash theorem]
  A strategic game $\langle N, (A_i),(u_i)\rangle$ with
  \begin{itemize}
  \item a finite number of players 
   \item a finite number of actions for each player
  \end{itemize}
has a Nash equilibrium in mixed strategies.
\end{theorem}
\begin{itemize}
\item proof uses Brouwer's fixed point theorem
\item to simplify notation, we prove the theorem for $2\times 2$ games
\end{itemize}
}

\fr{\ft{Proof of Nash theorem for $2\times2$ games}
  \begin{itemize}
  \item $2\times 2$ game:
    \begin{figure}[h]
      \centering
      \begin{game}{2}{2}
        \>L\>R\\
U    \>a,b\>c,d\\
D   \>e,f\>g,h
      \end{game}
    \end{figure}
  \item mixed strategy of P1: probability $\alpha\in[0,1]$  of playing U 
 \item mixed strategy of P2: probability $\beta\in[0,1]$  of playing L  
 \item idea of proof:
   \begin{itemize}
   \item define a function $f(\alpha,\beta)$ such that
     \begin{itemize}
     \item $f$ is continuous
     % \item if $(\alpha^*,\beta^*)$ is a Nash equilibrium, then $f(\alpha^*,\beta^*)=(\alpha^*,\beta^*)$
     \item if $f(\alpha^*,\beta^*)=(\alpha^*,\beta^*)$, then $(\alpha^*,\beta^*)$ is a Nash equilibrium of the game
     \end{itemize}
   \item use Brouwer's theorem to establish that $f$ has a fixed point
   \end{itemize}
  \end{itemize}

\framebreak
\begin{itemize}
\item $u_1(U,\beta)=\beta a+(1-\beta) c$ is the expected utility of P1 when playing U and P2 uses the mixed strategy $\beta$
\item $u_1(U,\beta)$ is linear and therefore continuous in $\beta$
\item define 
$$g(\alpha,\beta)=max\left\{0,\frac{\alpha+u_1(U,\beta)-u_1(D,\beta)}{1+\left|u_1(U,\beta)-u_1(D,\beta)\right|}\right\} $$

\begin{itemize}
\item $g(\alpha,\beta)$ is higher than $\alpha$ if U is the best response to $\beta$ and
  lower than $\alpha$ if D is best response 
\item $g$ is continuous because $u_1$ is continuous in $\beta$
\end{itemize}

\item define
$$h(\alpha,\beta)=max\left\{0,\frac{\beta+u_2(L,\alpha)-u_2(R,\alpha)}{1+\left|u_2(L,\alpha)-u_2(R,\alpha)\right|}\right\} $$
\begin{itemize}
\item $h(\alpha,\beta)$ is higher than $\beta$ if L is best response  to $\alpha$
  and lower than $\beta$ if R is best response
\item $h$ is
  continuous because $u_2$ is continuous in $\alpha$
\end{itemize}

\item let $f$ be
$$f(\alpha,\beta)=(g(\alpha,\beta),h(\alpha,\beta))$$
\end{itemize}

\begin{itemize}
\item if $f(\alpha^*,\beta^*)=(\alpha^*,\beta^*)$ then
  \begin{itemize}
  \item $g(\alpha^*,\beta^*)=\alpha^*$ $\Rightarrow$ $\alpha^*$ is best response to $\beta^*$
  \item $h(\alpha^*,\beta^*)=\beta^*$ $\Rightarrow$ $\beta^*$ is best response to $\alpha^*$
  \item $(\alpha^*,\beta^*)$ is Nash equilibrium
  \end{itemize}
every fixed point of $f$ is Nash equilibrium
\item  $f$ is continuous because $g$ and $h$ are continuous
\item Brouwer: $f$ has a fixed point
\end{itemize}
}

\fr{\ft{Existence revisited}
 \begin{ex}[revisited]
    \begin{itemize}
    \item 2 players ($N=\{P1,P2\}$)
    \item each player says a number ($A_i=\Re$)
    \item player that says the higher number wins \\(e.g. winner has payoff 1 while loser has payoff 0, if both say the same number each has payoff 1/2)
    \end{itemize}
  \end{ex}

  \begin{itemize}
  \item why does the Nash theorem not apply?
  \end{itemize}

Note that the propositions give sufficient conditions for NE, i.e. NE can exist even if the conditions are not satisfied:

\begin{ex}[war of attrition 1, OR ex. 18.4]
  \begin{itemize}
  \item 2 players dispute an object
  \item time is a continuous variable starting at 0 and running indefinitely
  \item actions: chose a time $a_i$ when to concede object
   \item player $i$ attaches value $v_i$ to object and loses one payoff unit for every time unit (until one of the players concedes)
  \end{itemize}
action space is not finite

% \begin{itemize}
% \item a action in this game is a concession time $a_i$
% \item can both players wait forever? No, then better to concede immediately
% \item can there be a NE where both players concede at the same time? No: better to wait $\varepsilon $ longer!
% \item can the loser have a strictly positive concession time? No: better to concede immediately
% \item when is  $a_i=0$ optimal? Only if $a_j\geq a_i$!
% \item equilibria: $a_i=0$ and $a_j\geq v_i$
% \end{itemize}
\end{ex}


}


\frame{\ft{Review Questions}
  \begin{itemize}
  \item Why are we interested in the question ``Does a Nash equilibrium exist?''?
  \item What is a ``fixed point''?
  \item Explain Brouwers fixed point theorem (1-dimensional case)!
  \item What does Nash's theorem state?
  \end{itemize}
reading: MSZ 5.3, (supplementary *reading: OR 2.4)
}

\fr{\ft{Exercises}
  \begin{enumerate}
  \item We look at a 2 player game in which $A_1=A_2=[0,1]$. Assume that P1 has a best response function that is strictly increasing and continuous. For all subquestions draw the best response functions in an $a_1,\, a_2$ diagram.
    \begin{enumerate}
    \item First, assume that P2 also has a best response function that is continuous. Does this guarantee the existence of a Nash equilibrium?
    \item Second, assume that P2 also has a best response function that is increasing (though not necessarily continuous).  Does this guarantee the existence of a Nash equilibrium?
   \item Third, assume that P2 has a best response function that is discontinuous at one point only. Does this guarantee the existence of a Nash equilibrium?
    \end{enumerate}
  \item *Use the Brouwer fixed point theorem to prove that a pure strategy Nash equilibrium exists in the following strategic game:
\begin{itemize}
\item $N=\{P1,P2\}$
\item $A_1=A_2=[0,1]$
\item $u_1(a_1,a_2)$ is differentiable and strictly concave in $a_1$. The partial derivative $\frac{\partial u_1}{\partial a_1}(a_1,a_2)$ is continuous in both variables.
\item $u_2(a_1,a_2)$ is differentiable and strictly concave in $a_2$.  The partial derivative $\frac{\partial u_2}{\partial a_2}(a_1,a_2)$ is continuous in both variables.
\end{itemize}
\emph{Hint: This exercise is an extension of the Cournot example we did in the lecture. The argument is similar.}
  \end{enumerate}
}

\end{document}

















to show items one by one it is easiest to insert a "\pause" command before the \item command

hyperlinks: \hyperlink{targetname}{\beamergotobutton{Text in button}} }
    to define a target insert the option [label=targetname] after the \frame command: \frame[label=targetname]{\ft{title}...}
    alternative the command \hypertarget{targetname}{} can be used
    To go back there exists a \beamerreturnbutton{text in button}
